% META: {"id": "1SPE-TRIGO-TD-FIL-ROUGE", "chapitre": "1SPE-TRIGONOMETRIE", "type_objet": "td", "capacites_codes": ["C1","C2"], "status": "generated"}

%---------------------------------------------------------------
% TD FIL ROUGE — Du radian aux coordonnées sur le cercle
% Couvre les capacites officielles C1 et C2 du programme 2026.
%---------------------------------------------------------------

\section*{TD fil rouge — Du radian aux coordonnées sur le cercle}

\subsection*{Étape 1 — Cercle trigonométrique et radian (C1)}

\begin{exercice}{1SPE-TRIGO-TD-FR-1}{1}{8}
\begin{enumerate}
  \item Convertir en radians : $150°$, $210°$, $315°$, $720°$.
  \item Convertir en degrés : $\dfrac{5\pi}{6}$, $\dfrac{7\pi}{4}$, $\dfrac{11\pi}{6}$, $3\pi$.
  \item Déterminer la mesure principale de chacun des angles suivants : $\dfrac{13\pi}{4}$, $-\dfrac{7\pi}{3}$, $\dfrac{25\pi}{6}$.
\end{enumerate}
\end{exercice}

% BEGIN-VERIFY
% from sympy import *
% # Conversions degres -> radians
% assert 150 * pi / 180 == 5*pi/6
% assert 210 * pi / 180 == 7*pi/6
% assert 315 * pi / 180 == 7*pi/4
% assert 720 * pi / 180 == 4*pi
% # Conversions radians -> degres
% assert 5*pi/6 * 180 / pi == 150
% assert 7*pi/4 * 180 / pi == 315
% assert 11*pi/6 * 180 / pi == 330
% assert 3*pi * 180 / pi == 540
% # Mesures principales
% # 13*pi/4 - 2*pi = 13*pi/4 - 8*pi/4 = 5*pi/4 -> encore > pi, 5*pi/4 - 2*pi = -3*pi/4
% assert Rational(13,4) - 2 == Rational(5,4)
% # mesure principale de 5*pi/4 est 5*pi/4 - 2*pi = -3*pi/4 (dans ]-pi;pi])
% END-VERIFY

\subsection*{Étape 2 — Valeurs remarquables et symetries (C2)}

\begin{exercice}{1SPE-TRIGO-TD-FR-2}{1}{10}
\begin{enumerate}
  \item Compléter le tableau des valeurs remarquables (sans calculatrice).
  \item Calculer $\cos\!\left(\dfrac{5\pi}{6}\right)$ et $\sin\!\left(\dfrac{5\pi}{6}\right)$ en utilisant la formule du supplementaire.
  \item Calculer $\cos\!\left(-\dfrac{3\pi}{4}\right)$ et $\sin\!\left(-\dfrac{3\pi}{4}\right)$.
  \item Vérifier que $\cos^2\!\left(\dfrac{5\pi}{6}\right) + \sin^2\!\left(\dfrac{5\pi}{6}\right) = 1$.
\end{enumerate}
\end{exercice}

% BEGIN-VERIFY
% from sympy import *
% assert cos(5*pi/6) == -sqrt(3)/2
% assert sin(5*pi/6) == Rational(1,2)
% assert cos(-3*pi/4) == -sqrt(2)/2
% assert sin(-3*pi/4) == -sqrt(2)/2
% assert trigsimp(cos(5*pi/6)**2 + sin(5*pi/6)**2) == 1
% END-VERIFY
